\section{Introduction}
\IEEEPARstart{T}{he} Internet of Things (IoT) has rapidly evolved from a conceptual framework into the foundational infrastructure of modern smart cities, industrial control systems, and healthcare networks. This ubiquitous connectivity, however, has exponentially expanded the global attack surface. Traditional Intrusion Detection Systems (IDS), historically deployed within the centralized confines of cloud architectures, face severe limitations in this new paradigm. Transmitting raw, high-velocity network telemetry from billions of distributed IoT endpoints to a centralized cloud introduces unacceptable latency, bandwidth saturation, and critical data privacy vulnerabilities.

\subsection{The Edge Security Paradox: Resources vs. Robustness}
To circumvent the limitations of cloud-centric security, the prevailing architectural shift advocates for \textit{Edge Computing}—deploying IDS immediately at the network periphery on localized IoT gateways, such as the Raspberry Pi. This localized approach promises near zero-latency threat mitigation. However, it introduces a profound paradox. The machine learning (ML) architectures that achieve state-of-the-art anomaly detection (e.g., Deep Neural Networks, complex Ensembles) are exceptionally computationally expensive. They demand extensive memory footprints and high-throughput processors that profoundly exceed the constraints of typical low-power, ARM-based edge devices. Consequently, deploying robust IDS at the edge frequently results in systemic bottlenecks, dropped packets, and catastrophic system failures during high-volume DDoS attacks.

\subsection{The Black-Box Trust Dilemma}
Compounding the hardware crisis is the fundamental opacity of modern AI. Deep Learning and complex ensemble models operate as "black-boxes," rendering their decision-making processes entirely invisible to human operators. In high-stakes enterprise and industrial environments, a binary alert (i.e., "Malicious" vs. "Benign") is operationally insufficient. Security Operations Center (SOC) analysts and network administrators require immediate, contextual data explaining \textit{why} a specific packet flow was flagged to execute a precise incident response. This lack of interpretability breeds a profound lack of trust, delaying critical mitigation efforts and rendering the automated system functionally opaque.

\subsection{Proposed Solution and Contributions}
To resolve the dichotomy between edge hardware constraints and the demand for trustworthy AI, this paper proposes \textbf{X-IDS}, a comprehensively Explainable, Resource-Efficient Intrusion Detection framework designed exclusively for edge deployment. 

Our framework leverages a lightweight machine learning core deployed directly onto a Raspberry Pi 4, ingesting real-time network traffic through a highly optimized, distributed \textbf{Apache Kafka} stream processing pipeline. Crucially, we integrate \textbf{Explainable AI (XAI)} into the core of the architecture by utilizing \textit{SHapley Additive exPlanations (SHAP)}. We diverge from traditional post-hoc analysis by employing SHAP as a dual-purpose engine: first, as a rigorous, mathematically grounded feature-selection mechanism to aggressively prune computational overhead before deployment; and second, as a real-time interpretation layer that assigns exact contributory weight to individual network features for every generated alert.

The primary contributions of this paper are defined as follows:
\begin{enumerate}
    \item \textbf{Trustworthy Edge Architecture:} We present the end-to-end design and empirical validation of an IDS operating on a Raspberry Pi, utilizing Apache Kafka to handle real-time traffic streams efficiently without localized hardware saturation.
    \item \textbf{Real-Time Interpretability via SHAP:} We integrate SHAP to crack the "black-box," providing security administrators with discrete, human-readable explanations (e.g., specific HTTP payload lengths or anomalous flow durations) driving every intrusion alert, thereby establishing operational trust.
    \item \textbf{XAI-Driven Resource Optimization:} We mathematically demonstrate that by utilizing SHAP values to identify and isolate only the most critical telemetry features, we drastically reduce the model's memory footprint and inference latency, achieving state-of-the-art accuracy within strict edge-computing hardware parameters.
\end{enumerate}

The remainder of this paper is structured as follows: Section II reviews current literature regarding Edge IDS and XAI. Section III formalizes the methodology and SHAP mechanics. Section IV details the implementation of the Raspberry Pi and Kafka architecture. Section V provides extensive empirical results and performance analysis, and Section VI concludes the research.
